\section*{Chapter \ref{ch:g11:ocean-ridges} --- Ocean Ridges and the Oceanic Lithosphere}

\begin{solution}{exo:g11:ocean-ridges:1}
\omterm{def:g11:ocean-ridges:peridotite}{Peridotite} is a dark, coarse-grained rock made mostly of olivine, with some
pyroxene; much of the mantle is made of it. \omterm{def:g11:ocean-ridges:gabbro}{Gabbro} is a dark, coarse-grained
igneous rock of plagioclase feldspar and pyroxene, the most abundant rock of
the deep \omterm{def:g11:ocean-ridges:oceanic-lithosphere}{oceanic crust}. \omterm{def:g11:ocean-ridges:gabbro}{Gabbro} belongs to the crust, \omterm{def:g11:ocean-ridges:peridotite}{peridotite} to the mantle.
\omterm{def:g11:ocean-ridges:gabbro}{Gabbro} crystallised slowly, deep under the sea floor, where the magma cooled
slowly: its crystals had time to grow. Basalt lava cooled quickly at the sea
floor, and its crystals stayed small.
\end{solution}

\begin{solution}{exo:g11:ocean-ridges:2}
From the sea floor down: sediments (none at the axis), basalt lavas (pillows
and sheet flows), sheeted dykes, \omterm{def:g11:ocean-ridges:gabbro}{gabbro}, then, below the \omterm{def:g11:seismic-structure:moho}{Moho}, the \omterm{def:g11:ocean-ridges:peridotite}{peridotite}
of the mantle. The hole crossed $1507/7100 \approx 0.21$, about a fifth of
the crust: the deep \omterm{def:g11:ocean-ridges:gabbro}{gabbros} and the \omterm{def:g11:seismic-structure:moho}{Moho} have never been reached by a drill.
\end{solution}

\begin{solution}{exo:g11:ocean-ridges:3}
$d(0) = \qty{2500}{m}$; $d(9) = 2500 + 350 \times 3 = \qty{3550}{m}$;
$d(25) = 2500 + 350 \times 5 = \qty{4250}{m}$.
\end{solution}

\begin{solution}{exo:g11:ocean-ridges:4}
A \omterm{def:g11:ocean-ridges:normal-fault}{normal fault} is an inclined fault along which the block above the fault
plane has moved down relative to the block below. At a slow ridge the plates
pull the cold, brittle \omterm{def:g11:lithosphere-asthenosphere:lithosphere}{lithosphere} apart: it is stretched, and it breaks. The
blocks above the faults slide down, toward the axis, and the axial valley is
the block that has dropped between them.
\end{solution}

\begin{solution}{exo:g11:ocean-ridges:5}
It does not boil because of the pressure: more than two kilometres of water
press on it, and under that pressure water stays liquid far above
\qty{100}{\celsius}. The black ``smoke'' is a cloud of tiny grains of iron
sulfides and other minerals, which precipitate when the hot vent water, rich
in dissolved metals, meets the cold sea water.
\end{solution}

\begin{solution}{exo:g11:ocean-ridges:6}
(1) About \qty{54}{km} down, at about \qty{1324}{\celsius}. (2) About
\qty{85}{km} down, at about $1400 + 0.45 \times 85 \approx
\qty{1438}{\celsius}$. (3) The hotter mantle starts to melt deeper, so a
longer column of rising rock melts before reaching the crust: it makes more
basalt, and a thicker crust.
\end{solution}

\begin{solution}{exo:g11:ocean-ridges:7}
(1) $1489 - 1345 = \qty{144}{\celsius}$ below its melting temperature.
(2) It rises $100 - 54 = \qty{46}{km}$; the gap closes by
$144/46 \approx \qty{3.1}{\celsius}$ per kilometre on average. (3) No: on the
way up the rock even cools a little, by about $0.45 \times 46 \approx
\qty{21}{\celsius}$. It is its melting temperature that falls, much faster,
as the pressure decreases.
\end{solution}

\begin{solution}{exo:g11:ocean-ridges:8}
The half rates are \qty{7.3}{cm} and \qty{1.65}{cm} a year, i.e.\
\qty{73}{km} and \qty{16.5}{km} per million years. To make \qty{1000}{km} of
floor on one side: $1000/73 \approx 13.7$ million years for the East Pacific
Rise, $1000/16.5 \approx 61$ million years for the Mid-Atlantic Ridge.
\end{solution}

\begin{solution}{exo:g11:ocean-ridges:9}
(1) \qty{5}{cm} is \qty{5e-5}{km}: $65000 \times 5 \times 10^{-5} =
\qty{3.25}{km^2}$ of new floor a year. (2) $3.25 \times 7.1 \approx
\qty{23}{km^3}$ of crust a year. (3) $3.6 \times 10^8 / 3.25 \approx 1.1
\times 10^8$ years, about 110 million years: at this pace the ridges could
make the whole ocean floor in about a hundred million years, a short time
compared with the age of the Earth. This is why no ocean floor is very old.
\end{solution}

\begin{solution}{exo:g11:ocean-ridges:10}
(1) The fault plane makes \ang{60} with the horizontal, and the drop is the
vertical side of a right triangle whose hypotenuse is the slip:
$1.5/\sin\ang{60} \approx \qty{1.73}{km}$. (2) Each fault moves the sides
apart horizontally by $1.5/\tan\ang{60} \approx \qty{0.87}{km}$; the two
faults together, by about \qty{1.73}{km}. (3) No. The \qty{25}{km} is the
distance between the faults that bound the dropped block, not the amount of
stretching, which is only about \qty{1.7}{km}; and the ridge keeps opening
by fresh dykes and lavas at the axis as well as by faulting.
\end{solution}

\begin{solution}{exo:g11:ocean-ridges:11}
$d(25) = 6400 - 3200\,e^{-25/62.8} \approx 6400 - 3200 \times 0.672 \approx
\qty{4251}{m}$; $d(100) = 6400 - 3200\,e^{-100/62.8} \approx 6400 - 3200
\times 0.203 \approx \qty{5749}{m}$. Between 25 and 100 million years the
floor sinks by about \qty{1500}{m}. The square-root law gives
$2500 + 350 \times 5 = \qty{4250}{m}$ at 25 million years: the two laws
agree, where both apply.
\end{solution}

\begin{solution}{exo:g11:ocean-ridges:12}
(1) $\alpha\,\Delta T\,L = 3.2 \times 10^{-5} \times 700 \times 100 =
\qty{2.24}{km}$. (2) $s = 3380 \times 2.24/(3380 - 1025) = 7571/2355 \approx
\qty{3.2}{km}$. (3) The exponential law allows $6400 - 2500 =
\qty{3900}{m}$ of sinking in all: cooling and contraction alone explain about
four fifths of it, with simple values for the temperature drop and the
thickness.
\end{solution}

\begin{solution}{exo:g11:ocean-ridges:13}
(1) $54 - 7 = \qty{47}{km}$ against $85 - 7 = \qty{78}{km}$: the hotter
column is $78/47 \approx 1.66$ times longer. (2) $7.1 \times 1.66 \approx
\qty{11.8}{km}$. (3) The measured \qty{10.3}{km} is thicker than normal, as
predicted, but less than the estimate: the deepest kilometres of the column,
just above the onset of melting, give less melt than the shallower ones, so
the assumption of the same melt per kilometre overestimates.
\end{solution}

\begin{solution}{exo:g11:ocean-ridges:14}
(1) The law gives \qty{3178}{m}, \qty{4160}{m}, \qty{4782}{m} and
\qty{5267}{m}; the measured medians are shallower by \qty{64}{m},
\qty{142}{m}, \qty{317}{m} and \qty{619}{m}. (2) The difference grows with
age. (3) The law gives the depth of the floor itself, corrected for its
sediments; the measured cells carry their sediments, which are thicker on
older floor, and include seamounts and volcanic plateaus, which are
shallower than the floor around them.
\end{solution}

\begin{solution}{exo:g11:ocean-ridges:15}
(1) $p = 1025 \times 9.8 \times 2500 \approx \qty{25.1}{MPa}$, plus
\qty{0.1}{MPa}: about \qty{25.2}{MPa}. (2) This is above the critical
pressure, \qty{22.064}{MPa}: the vent water cannot boil, however hot it is.
(3) The pressure reaches \qty{22.064}{MPa} at the depth
$(22.064 - 0.1) \times 10^6/(1025 \times 9.8) \approx \qty{2190}{m}$: every
vent deeper than about \qty{2200}{m} is safe from boiling.
\end{solution}

\begin{solution}{pb:g11:ocean-ridges:1}
\textbf{1.} $14.6/2 = \qty{7.3}{cm}$ a year, i.e.\ \qty{73}{km} per million
years.

\textbf{2.} At the depth of the ridge axis, $d(0) = \qty{2500}{m}$.

\textbf{3.} From the sea floor down: basalt lavas (pillows and sheet flows),
sheeted dykes and \omterm{def:g11:ocean-ridges:gabbro}{gabbro}, about \qty{7.1}{km} of crust in all; below the
\omterm{def:g11:seismic-structure:moho}{Moho}, the \omterm{def:g11:ocean-ridges:peridotite}{peridotite} of the mantle, to the base of the \omterm{def:g11:lithosphere-asthenosphere:lithosphere}{lithosphere}, which is
still very thin at the axis. Sediments will settle on it later.

\textbf{4.} As the plates move apart, the \omterm{def:g11:lithosphere-asthenosphere:lithosphere}{asthenosphere} rises beneath the
ridge. Its temperature falls only by about \qty{0.45}{\celsius} per
kilometre, while its melting temperature falls much faster as the pressure
decreases. About \qty{54}{km} down the two meet: the rock starts to melt, and
goes on melting as it rises, without any heat being added (\omterm{def:g11:ocean-ridges:partial-melting}{decompression
melting}).

\textbf{5.} Each million years the ridge makes $2 \times 73 = \qty{146}{km}$
of floor on its two sides; for each kilometre of ridge, the volume of crust is
$146 \times 7.1 \times 1 \approx \qty{1040}{km^3}$.

\textbf{6.} $73 \times t$: \qty{73}{km}, \qty{292}{km}, \qty{1168}{km} and
\qty{4672}{km}.

\textbf{7.} $2000/73 \approx 27$ million years.

\textbf{8.} Plates change their speed and direction over tens of millions of
years, as the bend of the Hawaiian--Emperor chain shows
(\cref{ch:g11:plate-motions}); the rate of today need not hold for 64
million years.

\textbf{9.} Away from the axis the crust is no longer above a magma lens: it
cools, and there is no longer enough heat below to drive the circulation and
heat the water.

\textbf{10.} West of the East Pacific Rise the piece travels on the Pacific
plate. Toward the west, the Pacific plate ends along deep trenches, where it
sinks beneath its neighbours: subduction zones (\cref{ch:g11:subduction-zones}).

\textbf{11.} $d(1) = \qty{2850}{m}$, $d(4) = 2500 + 700 = \qty{3200}{m}$,
$d(16) = 2500 + 1400 = \qty{3900}{m}$, $d(64) = 2500 + 2800 = \qty{5300}{m}$.

\textbf{12.} $3900 - 2500 = \qty{1400}{m}$ during the first 16 million years,
and $5300 - 3900 = \qty{1400}{m}$ between 16 and 64 million years: the same
sinking, in a time three times longer. The floor sinks more and more slowly
as it ages.

\textbf{13.} $350\sqrt{t} = 1500$ gives $\sqrt{t} \approx 4.29$, so
$t \approx 18.4$ million years.

\textbf{14.} $6400 - 3200\,e^{-64/62.8} \approx 6400 - 3200 \times 0.361
\approx \qty{5245}{m}$, about \qty{55}{m} less than the square-root law: the
two laws agree to about \qty{1}{\percent}.

\textbf{15.} The measured median is about \qty{650}{m} shallower than the
law. The law describes the floor corrected for the weight of its sediments;
the measured depths include the sediments, many tens of millions of years of
them on such floor, and the seamounts and volcanic plateaus, all of which make
the measured floor shallower.

\textbf{16.} $\alpha\,\Delta T = 3.2 \times 10^{-5} \times 700 = 0.0224$: the
mantle part is \qty{2.24}{\percent} denser,
$3380 \times 1.0224 \approx \qty{3456}{kg/m^3}$.

\textbf{17.} The mantle part is $100 - 7.1 = \qty{92.9}{km}$ thick. Mean
density: $(7.1 \times 2900 + 92.9 \times 3456)/100 \approx (20590 + 321060)/100
\approx \qty{3416}{kg/m^3}$.

\textbf{18.} $3416 > 3380$: the old \omterm{def:g11:lithosphere-asthenosphere:lithosphere}{lithosphere} as a whole is about
\qty{36}{kg/m^3}, about \qty{1}{\percent}, denser than the \omterm{def:g11:lithosphere-asthenosphere:lithosphere}{asthenosphere}
beneath it. The light crust is no longer enough to make up for its cold,
dense mantle part.

\textbf{19.} Being denser than the \omterm{def:g11:lithosphere-asthenosphere:lithosphere}{asthenosphere}, it can sink back into the
mantle at a trench: subduction, the subject of the next chapter.

\textbf{20.} At 64 million years the piece of floor lies at about
\qty{5300}{m}: about \textbf{5.3~km} down, \qty{2.8}{km} deeper than where it
was born.
\end{solution}
